% Generated by paper4/manuscript/build.py from paper4/paper/section_E9.md (e095dfdfe03d9840): edit the Markdown, not this file.
\subsection{How many: the coverage curve on whole cities (E9)}
E9 asks how many sensors it takes before instrumented neighbours start to inform the segments without one. On the whole networks of the three cities, the instrumented fraction \emph{f} takes four levels (10, 20, 40 and 60\%); the corridors of observed, validation-hidden and test-hidden segments are whole grid cells, and at every level the test-hidden segments fall in at least 8 cells. The arms are G0, G1 and G1 on a degree-preserving randomisation of the road graph (G1-rand(deg)), with five seeds: 180 trainings.

\textbf{No level of coverage makes propagation pay.} Propagation ``helps'' at a level if V(G1) has an interval above zero and a point estimate of at least 3 points. This happens in none of the 12 (city, \emph{f}) pairs. A 3-point advantage is excluded in 11 of 12; the exception is Melbourne at 60\%, −0.92 [−5.51; +3.25], so the pre-registered statement is ``undetermined''. The median half-width of V(G1) is 0.35 points, and the two intervals that exclude zero are unfavourable to the graph (Madrid 10\%, −0.39 [−0.94; −0.03]; London 20\%, −0.26 [−0.55; −0.00]).

\textbf{The curve is flat.} With a city term, a constant fits the 12 points better than a line (AIC −16.06 against −14.06) and better than a step (−12.25). Coverage changes the median distance from a hidden segment to the nearest instrumented one from 794 m at 10\% to 171 m at 60\% in Madrid, without moving V.

\begin{table}[!t]
\centering\footnotesize
\renewcommand{\arraystretch}{1.15}
\caption{How many: the coverage curve on whole cities (E9)}
\begin{tabularx}{\columnwidth}{>{\raggedright\arraybackslash\hsize=0.733\hsize}X>{\raggedright\arraybackslash\hsize=0.867\hsize}X>{\raggedright\arraybackslash\hsize=1.400\hsize}X}
\toprule
\textbf{City} & \textbf{wCE of G1 on test-hidden segments} & \textbf{Below ln 3 = 1.0986} \\
\midrule
Madrid & 1.075–1.087 & 1.1–2.2\% \\
London & 1.050–1.059 & 3.6–4.4\% \\
Melbourne & 1.081–1.134 & from −3.2\% to +1.6\% \\
\bottomrule
\end{tabularx}
\end{table}

On segments without a sensor the forecast stays close to the uniform predictor at every coverage, and red recall at 5\% false alarms is 0.05–0.08 in Madrid and London, against realised false-alarm rates of 0.04–0.06. The degree-preserving randomisation finds no geometry effect of 3 points or more (excluded in 11 of 12 pairs).

\begin{quote}
\textbf{At 10–60\% coverage, on whole cities, propagation from instrumented neighbours gives no measurable gain over the model without spatial exchange.}
\end{quote}

